\SourceLeaf{403}%
\[
\left.\begin{aligned}
\frac{r-q}{p^*-q}&=\frac{(a-b)(d-c)(ab-cd)^2}{(a-d)(b-c)(ad-bc)^2}.pp'\\
\frac{r-q}{p^*-r}&=\frac{(a-b)(d-c)(ab-cd)^2}{(a-c)(b-d)(ac-bd)^2}.pp'\\
\frac{p^*-r}{p^*-q}&=\frac{(a-c)(b-d)(ac-bd)^2}{(a-d)(b-c)(ad-bc)^2}
\end{aligned}\right\}\ \text{sind reelle Zahlen}
\]
oder \(p^*,q,r\) liegen auf einer Geraden senkrecht zu \(0p^*p\)
\par\medskip\Needspace{5\baselineskip}
\pdfbookmark[2]{V.}{carnot-zusatz-v}
\centerline{\textbf{V.}}\nobreak\smallskip
\{Einem gegebenen Kreis ein Vieleck einbeschreiben, dessen Seiten durch ebenso viele gegebene Punkte gehen. \textsc{Schumacher}.\}

Der Radius des Kreises sei \(r\), die Koordinaten der Eckpunkte des Vielecks seien
\[
\begin{array}{llll}
r\cos\varphi,&r\cos\varphi',&r\cos\varphi''&\text{etc.}\\
r\sin\varphi,&r\sin\varphi',&r\sin\varphi''&\text{etc.}
\end{array}
\]
und die Koordinaten der gegebenen Punkte, durch die die verlängerten Seiten des Vielecks gehen (die Seiten, die jeweils den ersten mit dem zweiten, den zweiten mit dem dritten Eckpunkt und so weiter verbinden), seien
\[
\begin{array}{llll}
a\cos A,&a'\cos A',&a''\cos A''&\text{etc.}\\
a\sin A,&a'\sin A',&a''\sin A''&\text{etc.}
\end{array}
\]
\par
Liegen drei Punkte mit den Koordinaten \(x,y;\ x',y';\ x'',y''\) auf einer Geraden, so gilt die Bedingungsgleichung
\[
xy'+x'y''+x''y-x'y-x''y'-xy''=0
\]
Daraus folgt
\[
\left.\begin{aligned}
rr\cos\varphi\sin\varphi'+ra\cos\varphi'\sin A+ar\cos A\sin\varphi\\
-rr\cos\varphi'\sin\varphi-ar\cos A\sin\varphi'-ra\cos\varphi\sin A
\end{aligned}\right\}=0
\]
oder
\[
r\sin(\varphi'-\varphi)-a\sin(\varphi'-A)+a\sin(\varphi-A)=0
\]
oder
\[
r\cos\tfrac12(\varphi'-\varphi)=a\cos\tfrac12(\varphi'+\varphi-2A)
\]
